\documentclass[aspectratio=169,11pt]{beamer}

\usepackage{uski-beamer}

% Metadata
\title{USki -- Moderne Flashcards}
\subtitle{Spaced Repetition \& AI-Chat}
\author{Leon Huber}
\date{\today}
\institute{HTBLuVA Salzburg}

\begin{document}

% Title slide
\begin{frame}[plain]
  \titlepage
\end{frame}

% Outline
\begin{frame}{Agenda}
  \tableofcontents
\end{frame}

% ---------------------------------------------------------------
\section{Projektübersicht}

\begin{frame}{Was ist USki?}
  \begin{columns}[T]
    \begin{column}{0.55\textwidth}
      \textbf{USki} ist eine containerisierte Flashcard-App mit:
      \begin{itemize}
        \item \textbf{FSRS}-Algorithmus für Spaced Repetition
        \item \textbf{NotebookLM}-artiger KI-Chat
        \item Rich-Text Editor für Karten
        \item Sicheres RBAC-Sharing-System
      \end{itemize}
      \vspace{8pt}
      \textbf{Kernziele:}
      \begin{itemize}
        \item Effizientes, nachhaltiges Lernen
        \item Nahtlose AI-Unterstützung
        \item Moderne, minimalistische UX
      \end{itemize}
    \end{column}
    \begin{column}{0.40\textwidth}
      \begin{block}{Technologie-Stack}
        \textbf{Frontend:} React, Vite, TS, Tailwind\\[4pt]
        \textbf{Backend:} Python, FastAPI, Pydantic\\[4pt]
        \textbf{Cloud:} Supabase, Resend\\[4pt]
        \textbf{KI:} Google Gemini
      \end{block}
    \end{column}
  \end{columns}
\end{frame}

% ---------------------------------------------------------------
\section{Architektur}

\begin{frame}{Systemarchitektur}
  \begin{center}
    \resizebox{0.9\textwidth}{!}{%
    \begin{tikzpicture}
      % Nodes
      \node[accentbox] (fe) {Frontend\\ React / Vite};
      \node[accentbox, right=of fe] (be) {Backend\\ FastAPI};
      \node[box, below=1.5cm of be] (supa) {Supabase Cloud\\ Auth, DB, Storage};
      \node[box, right=of be] (gem) {Google Gemini\\ API};
      \node[box, below=1.5cm of fe] (dock) {Docker Compose\\ Frontend + Backend};
      % Arrows
      \draw[arrow, <->] (fe) -- (be) node[midway, above, mylabel] {REST API};
      \draw[arrow, <->] (be) -- (supa) node[midway, right, mylabel] {Auth / DB};
      \draw[arrow, <->] (be) -- (gem) node[midway, above, mylabel] {KI};
      \draw[arrow, dashed, gray] (fe) -- (dock);
      \draw[arrow, dashed, gray] (be) -- (dock);
    \end{tikzpicture}
    }
  \end{center}
\end{frame}

% ---------------------------------------------------------------
\section{Authentifizierung}

\begin{frame}{Passwortloser Login-Flow}
  \textbf{Keine Passwörter. Kein Reset. Nur E-Mail-OTP.}
  \vspace{12pt}
  \begin{center}
    \resizebox{0.9\textwidth}{!}{%
    \begin{tikzpicture}
      % Nodes
      \node[box] (user) {User};
      \node[box, right=2cm of user] (supa) {Supabase Auth};
      \node[box, right=2cm of supa] (resend) {Resend SMTP};
      \node[box, below=1.8cm of supa] (api) {FastAPI};
      % Flow
      \draw[arrow, ->] (user) -- (supa) node[midway, above, mylabel] {E-Mail};
      \draw[arrow, ->] (supa) -- (resend) node[midway, above, mylabel] {OTP-Code};
      \draw[arrow, ->] (resend) to[bend left=30] node[midway, below, mylabel, xshift=-0.5cm] {6-stell. Code} (user);
      \draw[arrow, ->] (user) -- (api) node[midway, left, mylabel] {JWT Token};
      \draw[arrow, ->] (api) -- (supa) node[midway, right, mylabel] {Verify};
    \end{tikzpicture}
    }
  \end{center}
\end{frame}

% ---------------------------------------------------------------
\section{API-Design}

% [fragile] is required for lstlisting inside a frame
\begin{frame}[fragile]{FastAPI Router}
  \begin{columns}[T]
    \begin{column}{0.50\textwidth}
      \lstset{language=Python}
      \begin{lstlisting}
# backend/src/uski/api/router.py
from fastapi import APIRouter

router = APIRouter(prefix="/api/v1")

@router.get("/health")
async def health_check():
    return {"status": "ok"}
      \end{lstlisting}
    \end{column}
    \begin{column}{0.46\textwidth}
      \textbf{Struktur:}
      \begin{itemize}
        \item Health-Check
        \item Auth-Router (JWT)
        \item Decks, Cards, Chat
        \item File-Upload
      \end{itemize}
      \vspace{6pt}
      \begin{block}{Schema-Validierung}
        Alle Request/Response-Modelle mit \textbf{Pydantic v2}
      \end{block}
    \end{column}
  \end{columns}
\end{frame}

% ---------------------------------------------------------------
\section{KI-Integration}

\begin{frame}{RAG-Pipeline}
  \begin{columns}[T]
    \begin{column}{0.50\textwidth}
      \textbf{Retrieval Augmented Generation}
      \begin{enumerate}
        \item User fragt KI-Chat
        \item Query wird zu Embedding
        \item Relevante \texttt{document\_chunks} aus Supabase (pgvector)
        \item Kontext + Prompt $\to$ Gemini
        \item Antwort + Sources an Nutzer
      \end{enumerate}
    \end{column}
    \begin{column}{0.46\textwidth}
      \begin{block}{Embeddings}
        Google Gemini 1.5\\
        \texttt{text-embedding}
      \end{block}
      \vspace{6pt}
      \begin{block}{Model}
        Gemini 1.5 Flash\\
        Text + Vision
      \end{block}
    \end{column}
  \end{columns}
\end{frame}

% ---------------------------------------------------------------
\section{Zusammenfassung}

\begin{frame}{Ausblick}
  \begin{columns}[T]
    \begin{column}{0.48\textwidth}
      \begin{block}{In Arbeit}
        \begin{itemize}
          \item Docker-Compose finalisieren
          \item Frontend-Views implementieren
          \item Auth-Flow vollständig testen
        \end{itemize}
      \end{block}
    \end{column}
    \begin{column}{0.48\textwidth}
      \begin{block}{Zukunft}
        \begin{itemize}
          \item Mobile App (React Native)
          \item Offline-Support (PWA)
          \item Gamification
        \end{itemize}
      \end{block}
    \end{column}
  \end{columns}
\end{frame}

\begin{frame}[plain]
  \centering
  \vfill
  {\Large\textbf{Vielen Dank!}}

  \vspace{12pt}
  {\color{graytext} Fragen?}
  \vfill
\end{frame}

\end{document}
